\documentclass[11pt,a4paper]{report}

\usepackage[utf8]{inputenc}
\usepackage[T1]{fontenc}
\usepackage[french]{babel}
\usepackage[margin=2.1cm]{geometry}
\usepackage{xurl}
\usepackage{hyperref}
\usepackage{booktabs}
\usepackage{xcolor}
\usepackage{longtable}
\usepackage{enumitem}
\usepackage{fancyhdr}
\usepackage{titlesec}
\usepackage{array}
\usepackage{listings}

\definecolor{copifiobsidian}{HTML}{050505}
\definecolor{copifiemerald}{HTML}{10B981}
\definecolor{copifisky}{HTML}{0EA5E9}
\definecolor{copifizinc500}{HTML}{71717A}
\definecolor{copifiwarn}{HTML}{F59E0B}

\hypersetup{
    colorlinks=true,
    linkcolor=copifiemerald,
    urlcolor=copifisky,
    pdftitle={Compte rendu — Dashboards Copifi},
    pdfauthor={Équipe Copifi}
}

\pagestyle{fancy}
\fancyhf{}
\fancyhead[L]{\small\textit{Dashboards Copifi — inventaire détaillé}}
\fancyhead[R]{\small\thepage}
\renewcommand{\headrulewidth}{0.4pt}

\titleformat{\chapter}[display]
{\normalfont\huge\bfseries\color{copifiobsidian}}{\chaptertitlename\ \thechapter}{14pt}{\Huge}

\titleformat{\section}
{\normalfont\Large\bfseries\color{copifiemerald}}
{\thesection}{1em}{}

\newcommand{\filepath}[1]{\texttt{\small #1}}
\newcommand{\route}[1]{\texttt{\small #1}}

\begin{document}

\begin{titlepage}
    \centering
    \vspace*{1cm}
    {\Huge\bfseries\color{copifiobsidian} Compte rendu\par}
    \vspace{0.5cm}
    {\LARGE\color{copifiemerald} Contenu des dashboards Copifi\par}
    \vspace{0.35cm}
    {\large Inventaire détaillé de ce qui est affiché, calculé et accessible\par}
    \vspace{1.2cm}
    \rule{0.82\textwidth}{1pt}
    \vspace{0.9cm}
    \begin{tabular}{>{\bfseries}r p{9.4cm}}
        Produit & Copifi \\
        Périmètre & Zone connectée après abonnement actif \\
        Fichiers clés & \filepath{Dashboard.jsx}, \filepath{Facturation/Dashboard.jsx}, \filepath{AppDashboardLayout.jsx} \\
        Route principale pilotage & \route{/dashboard} \\
        Route dashboard facturation & \route{/facturation} \\
        Version & Juillet 2026 \\
    \end{tabular}
    \vfill
    \small\itshape
    Document basé sur le code source actuel (React/Inertia + Laravel).
\end{titlepage}

\tableofcontents
\cleardoublepage

% ============================================================
\chapter{Vue d'ensemble}
% ============================================================

Copifi expose \textbf{plusieurs écrans « dashboard »} dans l'espace connecté :

\begin{center}
\begin{tabular}{@{}p{3.2cm}p{2.5cm}p{7.8cm}@{}}
\toprule
\textbf{Écran} & \textbf{Route} & \textbf{Rôle} \\
\midrule
Tableau de bord mensuel & \route{/dashboard} & Pilotage financier (CA, marge, CAC, LTV, IA) \\
Vue facturation & \route{/facturation} & KPIs facturation, graphiques, activité documents \\
Administration & \route{/admin} & Gestion utilisateurs (rôle admin uniquement) \\
Copilote IA & \route{/copilote} & Chat conversationnel (hors dashboard strict) \\
Saisie mensuelle & \route{/saisie-mensuelle} & Formulaire alimentant le dashboard pilotage \\
\bottomrule
\end{tabular}
\end{center}

Tous les dashboards utilisateur passent par le \textbf{même shell visuel} :
\filepath{resources/js/Layouts/AppDashboardLayout.jsx} (sidebar Copifi, header, menu mobile).

\textbf{Accès :} middleware \texttt{auth}, \texttt{verified}, \texttt{active.subscription} (abonnement Stripe \texttt{active|trialing}, compte non suspendu).

% ============================================================
\chapter{Shell commun — AppDashboardLayout}
% ============================================================

\section{Sidebar (navigation)}

Pour un \textbf{utilisateur standard}, la sidebar contient 3 sections :

\subsection{Section « Pilotage »}
\begin{itemize}[leftmargin=1.5em]
    \item Tableau de bord $\rightarrow$ \route{/dashboard}
    \item Copilote IA $\rightarrow$ \route{/copilote}
    \item Saisie mensuelle $\rightarrow$ \route{/saisie-mensuelle}
\end{itemize}

\subsection{Section « Facturation »}
\begin{itemize}[leftmargin=1.5em]
    \item Vue facturation $\rightarrow$ \route{/facturation}
    \item Liste des factures $\rightarrow$ \route{/factures}
    \item Création de facture $\rightarrow$ \route{/factures/nouveau}
    \item Liste des devis $\rightarrow$ \route{/devis}
    \item Création de devis $\rightarrow$ \route{/devis/nouveau}
    \item Factures fournisseurs $\rightarrow$ \route{/achats/factures-fournisseurs}
    \item Gestion des clients $\rightarrow$ \route{/clients}
    \item Catalogue produits $\rightarrow$ \route{/catalogue}
    \item Suivi des paiements $\rightarrow$ \route{/paiements}
    \item Paramètres $\rightarrow$ \route{/parametres}
\end{itemize}

\subsection{Section « Compte »}
\begin{itemize}[leftmargin=1.5em]
    \item Profil $\rightarrow$ \route{/profile}
\end{itemize}

\section{Header (barre supérieure)}

\begin{itemize}[leftmargin=1.5em]
    \item \textbf{Titre de page} dynamique (ex.\ « Tableau de bord mensuel »)
    \item \textbf{Badge période} ou statut (ex.\ mois courant, « Abonnement : actif »)
    \item Indicateur \textbf{Abonnement : actif} (pastille verte)
    \item Bouton \textbf{Notifications} (cloche + point rouge — \textcolor{copifiwarn}{non branché fonctionnellement})
    \item Bloc utilisateur : nom, rôle, initiale, bouton \textbf{Déconnexion}
\end{itemize}

\section{Responsive}

\begin{itemize}[leftmargin=1.5em]
    \item Sidebar fixe \texttt{w-64} visible à partir de \texttt{lg} (1024\,px)
    \item Menu \textbf{hamburger} + drawer mobile en dessous de \texttt{lg}
\end{itemize}

% ============================================================
\chapter{Dashboard pilotage — \route{/dashboard}}
% ============================================================

\section{Identité de l'écran}

\begin{itemize}[leftmargin=1.5em]
    \item \textbf{Titre affiché :} « Tableau de bord mensuel »
    \item \textbf{Badge header :} mois du dernier enregistrement financier (ex.\ « 2026-07 ») ou « Aucune période »
    \item \textbf{Page React :} \filepath{resources/js/Pages/Dashboard.jsx}
    \item \textbf{Contrôleur :} \filepath{DashboardController.php}
    \item \textbf{Services :} \filepath{FinancialService.php}, \filepath{AiInsightService.php}
    \item \textbf{Source de données :} table \texttt{financial\_records} (saisie via \route{/saisie-mensuelle})
\end{itemize}

\section{Bandeau admin (conditionnel)}

Si un administrateur consulte le dashboard \textbf{d'un autre utilisateur} (\route{/admin/users/\{user\}/dashboard}) :

\begin{itemize}[leftmargin=1.5em]
    \item Encart « Vue administrateur »
    \item Nom et email de l'utilisateur consulté
    \item Statut Stripe affiché
\end{itemize}

\section{Grille KPI — 4 cartes (section \#kpi-grid)}

Grille responsive : \texttt{1 col} $\rightarrow$ \texttt{md:2} $\rightarrow$ \texttt{xl:4}.

\begin{longtable}{@{}p{3.5cm}p{10.8cm}@{}}
\toprule
\textbf{Carte} & \textbf{Contenu affiché} \\
\midrule
\endhead
Chiffre d'affaires &
Montant du \textbf{dernier mois saisi} (format compact k\,\€). Sparkline 7 points sur l'historique des revenus mensuels. \\
Marge nette &
Pourcentage marge = (CA $-$ charges) / CA. Sous-texte : montant net en €. Sparkline des marges \% par mois. \\
CAC &
Coût d'acquisition client = \textbf{budget marketing / nombre clients} du mois. Affiche « N/A » si \texttt{clients\_count = 0}. \\
LTV &
Lifetime value simplifiée = \textbf{CA / nombre clients} du mois. « N/A » si pas de clients. Sparkline. \\
\bottomrule
\end{longtable}

\subsection{Formules backend (\filepath{FinancialService.php})}

\begin{itemize}[leftmargin=1.5em]
    \item \texttt{marge\_nette} = \texttt{revenue} $-$ \texttt{charges}
    \item \texttt{cac} = \texttt{marketing\_budget} / \texttt{clients\_count}
    \item \texttt{ltv} = \texttt{revenue} / \texttt{clients\_count}
\end{itemize}

\section{Zone centrale — graphique principal (\#main-chart-section)}

\textbf{Titre :} « Revenus vs Charges Totales »

\begin{itemize}[leftmargin=1.5em]
    \item Graphique \textbf{AreaChart} (Recharts) : courbes Revenus (bleu néon) et Charges (orange)
    \item Axe X : mois des enregistrements \texttt{financial\_records}
    \item Axe Y : montants en k\,\€
    \item Tooltip au survol avec détail par série
    \item État vide : message « Aucune donnée financière » + invitation à saisir via \route{/saisie-mensuelle}
    \item Boutons \textbf{12M} / \textbf{YTD} : \textcolor{copifiwarn}{affichés mais non fonctionnels} (pas de filtre branché)
\end{itemize}

\section{Tableau « Flux Récents » (\#recent-transactions)}

\begin{itemize}[leftmargin=1.5em]
    \item Affiche les \textbf{3 derniers mois} de l'historique financier
    \item Colonnes : Date (mois), Description (« Période financière »), Catégorie (badge « Revenu »), Montant (CA du mois)
    \item Lien « Voir tout » : \textcolor{copifiwarn}{présent sans navigation}
\end{itemize}

\section{Colonne droite — Score de santé (\#health-score)}

\textbf{Score /100} calculé côté frontend à partir des KPI :

\begin{itemize}[leftmargin=1.5em]
    \item Base 50 points
    \item Marge nette $>$ 0 : +20 ; $<$ 0 : $-$25
    \item Ratio LTV/CAC $>$ 3 : +20 ; entre 1 et 3 : +8 ; $<$ 1 : $-$18
    \item Charges / CA $\leq$ 50\,\% : +10 ; $\leq$ 70\,\% : +5 ; sinon $-$10
    \item Plafonné entre 0 et 100
    \item Visualisation : \textbf{donut chart} (anneau vert + reste gris)
\end{itemize}

\section{Panneau Alertes (\#alerts-panel)}

Alertes générées par \filepath{FinancialService::buildAlert()} + règles UI :

\begin{longtable}{@{}p{2.5cm}p{11.8cm}@{}}
\toprule
\textbf{Niveau} & \textbf{Déclencheur} \\
\midrule
\endhead
critique & Marge négative ; ou LTV $<$ CAC \\
attention & Charges $>$ 70\,\% du CA \\
sain & Marge $>$ 0 et LTV/CAC $>$ 3 \\
neutral & Données CAC/LTV incomplètes ; ou aucune alerte \\
\bottomrule
\end{longtable}

Chaque alerte : titre coloré + message explicatif.

\section{Bloc Analyse IA}

\textbf{Titre :} « Lecture financière du mois »

\begin{itemize}[leftmargin=1.5em]
    \item Texte généré par \textbf{Groq} via \filepath{AiInsightService} (stocké en BDD \texttt{ai\_insights})
    \item Bouton « Afficher / Masquer le diagnostic »
    \item Statuts : analyse disponible, « en attente », ou message si pas de saisie mensuelle
    \item Indication si texte \textbf{corrigé par un admin}
    \item \textbf{Mode admin} (vue utilisateur) : formulaire d'édition du texte IA + enregistrement (\route{PATCH /admin/ai-insights/\{id\}})
\end{itemize}

\section{Ce que ce dashboard ne contient PAS}

\begin{itemize}[leftmargin=1.5em]
    \item Pas de données factures/devis en direct (uniquement saisie manuelle mensuelle)
    \item Pas de widget chat copilote intégré sur cette page (chat sur \route{/copilote})
    \item Pas d'export PDF/Excel du dashboard pilotage
\end{itemize}

% ============================================================
\chapter{Dashboard facturation — \route{/facturation}}
% ============================================================

\section{Identité de l'écran}

\begin{itemize}[leftmargin=1.5em]
    \item \textbf{Titre :} « Vue d'ensemble »
    \item \textbf{Description :} « Gérez vos finances, factures et devis en temps réel. »
    \item \textbf{Page :} \filepath{Facturation/Dashboard.jsx} $\rightarrow$ composant \filepath{DashboardHome.jsx}
    \item \textbf{Layout :} \filepath{FacturationLayout.jsx} (en-tête + actions)
    \item \textbf{Contrôleur :} \filepath{FacturationDashboardController.php}
    \item \textbf{Source :} documents (\texttt{factures}, \texttt{devis}, \texttt{avoirs}), tiers, paiements
\end{itemize}

\section{Actions en-tête (\filepath{DashboardHeaderActions})}

\begin{itemize}[leftmargin=1.5em]
    \item Bouton \textbf{Rapport} (icône téléchargement) — \textcolor{copifiwarn}{non branché}
    \item Bouton \textbf{Créer Facture} $\rightarrow$ \route{/factures/nouveau}
\end{itemize}

\section{4 cartes KPI (première ligne)}

\begin{longtable}{@{}p{3.5cm}p{10.8cm}@{}}
\toprule
\textbf{KPI} & \textbf{Données affichées} \\
\midrule
\endhead
CA Encaissé &
Somme TTC en EUR des factures statut \texttt{paid}. Tendance \% vs mois précédent. Sparkline 7 mois (k\,\€). \\
Devis en attente &
Montant TTC des devis statut \texttt{sent} + nombre. Sparkline 7 mois. \\
Factures en retard &
Nombre de factures \texttt{sent} dont \texttt{due\_date} $<$ aujourd'hui. Sparkline du nombre sur 7 mois. \\
Nouveaux clients &
Clients créés ce mois-ci. Tendance \% vs mois précédent. Sparkline 7 mois. \\
\bottomrule
\end{longtable}

\section{Carte Analyse IA facturation (\filepath{FinancialAnalysisCard.jsx})}

\begin{itemize}[leftmargin=1.5em]
    \item Bouton « Générer l'analyse »
    \item Appel API \route{POST /facturation/analyze} $\rightarrow$ \filepath{FinancialAnalysisService} (Groq)
    \item Données envoyées (anonymisées) : CA encaissé mois, tendance, factures en retard, devis en attente, nouveaux clients
    \item Réponse affichée en \textbf{Markdown} (titres, listes, gras)
    \item États : chargement (skeleton), erreur session/503, succès
\end{itemize}

\section{Graphiques (deuxième ligne)}

\subsection{Évolution du chiffre d'affaires (2/3 largeur)}

\begin{itemize}[leftmargin=1.5em]
    \item AreaChart \textbf{6 derniers mois}
    \item Valeur = factures émises (hors annulées) $-$ avoirs envoyés, en k\,\€
    \item Bouton « 6 Derniers Mois » — \textcolor{copifiwarn}{décoratif, pas de changement de période}
\end{itemize}

\subsection{Répartition des factures (donut, 1/3)}

\begin{itemize}[leftmargin=1.5em]
    \item Camembert : \textbf{Payées} (vert), \textbf{En attente} (ambre), \textbf{En retard} (rouge)
    \item Centre : nombre total de factures actives
    \item Pourcentages calculés côté backend
\end{itemize}

\section{Activité récente + Raccourcis (troisième ligne)}

\subsection{Tableau Activité Récente (7 dernières lignes)}

Documents triés par \texttt{updated\_at} décroissant :

\begin{center}
\begin{tabular}{@{}ll@{}}
\toprule
\textbf{Colonne} & \textbf{Contenu} \\
\midrule
Client / Réf & Avatar initiales, nom client, référence document \\
Type & Facture, Devis ou Avoir \\
Montant & TTC devise document \\
Statut & Badge coloré (Payée, Envoyé, En retard…) \\
Date & Relatif (Aujourd'hui, Hier, Il y a Xj…) \\
\bottomrule
\end{tabular}
\end{center}

Lien « Voir tout » $\rightarrow$ \route{/factures}.

\subsection{Panneau Raccourcis}

Trois cartes cliquables :
\begin{itemize}[leftmargin=1.5em]
    \item Nouvelle Facture $\rightarrow$ \route{/factures/nouveau}
    \item Nouveau Devis $\rightarrow$ \route{/devis/nouveau}
    \item Ajouter Client $\rightarrow$ \route{/clients}
\end{itemize}

% ============================================================
\chapter{Dashboard admin — \route{/admin}}
% =============================================================

\section{Accès}

Réservé aux utilisateurs \texttt{role = admin}. Sidebar réduite : Admin + Profil.

\section{Contenu (\filepath{Admin/Dashboard.jsx})}

\subsection{4 cartes statistiques}
\begin{itemize}[leftmargin=1.5em]
    \item Utilisateurs (total inscrits)
    \item Abonnements valides
    \item MRR (revenu mensuel récurrent)
    \item CA total clients (somme des \texttt{financial\_records.revenue})
\end{itemize}

\subsection{Tableau utilisateurs}

Colonnes typiques : nom, email, statut Stripe (badge coloré), rôle, date inscription, actions :
\begin{itemize}[leftmargin=1.5em]
    \item Voir dashboard utilisateur
    \item Suspendre / Réactiver
\end{itemize}

\subsection{Journal d'audit}

Liste des dernières actions admin (\texttt{admin\_audit\_logs}) : suspension, restauration.

% ============================================================
\chapter{Pages liées au pilotage (hors dashboard strict)}
% ============================================================

\section{Saisie mensuelle — \route{/saisie-mensuelle}}

Formulaire alimentant le dashboard \route{/dashboard} :

\begin{itemize}[leftmargin=1.5em]
    \item Mois (YYYY-MM)
    \item Chiffre d'affaires (\texttt{revenue})
    \item Charges (\texttt{charges})
    \item Budget marketing (\texttt{marketing\_budget})
    \item Nombre de clients (\texttt{clients\_count})
    \item Envoi : \route{POST /financial-records}
\end{itemize}

\section{Copilote IA — \route{/copilote}}

Page dédiée avec \filepath{DashboardChatWidget.jsx} :
\begin{itemize}[leftmargin=1.5em]
    \item Chat Groq avec contexte financier + facturation
    \item API : \route{POST /copilote/chat} ou \route{/dashboard/chat}
    \item Limite : 20 messages/minute
\end{itemize}

% ============================================================
\chapter{Synthèse des sources de données}
% ============================================================

\begin{center}
\begin{tabular}{@{}p{4.5cm}p{4cm}p{5.8cm}@{}}
\toprule
\textbf{Dashboard} & \textbf{Table / source} & \textbf{Mise à jour} \\
\midrule
Pilotage KPI & \texttt{financial\_records} & Saisie manuelle mensuelle \\
Pilotage graphiques & \texttt{financial\_records} & Idem \\
Pilotage alertes & Calcul \filepath{FinancialService} & Temps réel à l'affichage \\
Pilotage IA & \texttt{ai\_insights} + Groq & Génération à la visite \\
Facturation KPI & \texttt{documents}, \texttt{tiers} & Temps réel documents \\
Facturation activité & \texttt{documents} + relations & Temps réel \\
Facturation analyse IA & Agrégats anonymisés + Groq & Sur clic utilisateur \\
Admin stats & \texttt{users}, Stripe, \texttt{financial\_records} & Temps réel \\
\bottomrule
\end{tabular}
\end{center}

% ============================================================
\chapter{Éléments UI présents mais non finalisés}
% ============================================================

\begin{itemize}[leftmargin=1.5em]
    \item Boutons \textbf{12M / YTD} sur graphique dashboard pilotage
    \item Lien \textbf{« Voir tout »} flux récents (pilotage)
    \item Bouton \textbf{Rapport} (dashboard facturation)
    \item Bouton \textbf{Notifications} (header) — pas de panneau notifications
    \item Filtre période graphique CA facturation (bouton « 6 Derniers Mois »)
\end{itemize}

\vfill
\begin{center}
\small\color{copifizinc500}
Compte rendu dashboards Copifi — Juillet 2026\\
\filepath{docs/compte-rendu-dashboard-copifi.tex}
\end{center}

\end{document}
